% ============================================================================
% Chapitre 2 — État de l'art
% ============================================================================
\chapter{État de l'art}

\section*{Introduction}

Ce chapitre examine les solutions de supervision existantes, les techniques de
détection d'anomalies disponibles, et les avancées récentes en matière de modèles
de langage appliqués aux opérations informatiques. Cette analyse comparative est
indispensable pour justifier les choix architecturaux de \caimp{} et démontrer
la valeur ajoutée de l'approche retenue.

\section{Solutions de supervision commerciales}

\subsection{Datadog}

Datadog est une plateforme \textsc{saas} de monitoring fondée en 2010, devenue
l'une des références du marché avec plus de 27~000 clients en 2023.
Son architecture repose sur un agent léger (\texttt{datadog-agent}) installé
sur chaque hôte, qui collecte métriques, traces et journaux et les envoie
vers les datacenters Datadog.

\textbf{Points forts :} intégrations natives avec plus de 700 technologies,
tableau de bord personnalisable, fonctionnalités d'apprentissage automatique
(\textit{Watchdog}) pour la détection d'anomalies, corrélation automatique
des signaux, \textsc{apm} (\textit{Application Performance Monitoring}).

\textbf{Limites :} tarification à l'hôte (15~\$/mois par hôte pour l'offre
Infrastructure, 31~\$/mois avec APM) rédhibitoire pour les petites équipes ;
données hébergées aux États-Unis, posant des questions de conformité \textsc{rgpd} ;
aucun déploiement \textit{on-premise} disponible ; les explications générées
par l'IA restent superficielles et ne s'appuient pas sur un \textsc{llm} generatif.

\subsection{Grafana + Prometheus}

La combinaison Grafana / Prometheus est devenue le standard \textit{de facto}
de l'observabilité open source. Prometheus \cite{timescaledb2017} est un système
de surveillance pull-based avec une base de données temporelles propriétaire
(TSDB) et un langage de requête dédié (PromQL). Grafana est un outil de
visualisation connectant des dizaines de sources de données.

\textbf{Points forts :} entièrement open source, communauté active, excellent
pour les architectures Kubernetes, grande flexibilité des alertes
(Alertmanager), intégrations étendues.

\textbf{Limites :} absence totale d'IA native ; TSDB de Prometheus non conçue
pour les requêtes analytiques complexes (pas de fenêtres glissantes, pas
d'agrégats temporels avancés) ; nécessite une expertise significative pour la
configuration ; génère facilement de la fatigue d'alertes sans déduplication
intelligente.

\subsection{New Relic}

New Relic est une plateforme \textsc{saas} centrée sur la performance applicative
(\textsc{apm}) avec des capacités d'infrastructure monitoring. Son architecture
Telemetry Data Platform ingère plusieurs types de signaux dans un \textit{data lake}
unifié. La fonctionnalité \textit{New Relic AI} (NR AI) propose une détection
d'anomalies basée sur des algorithmes de machine learning propriétaires.

\textbf{Limites :} modèle de facturation à la consommation de données imprévisible
pour les charges variables ; l'IA reste un moteur de classification, sans
explication en langage naturel.

\subsection{Dynatrace}

Dynatrace se distingue par son moteur \textit{Davis AI}, un système d'analyse
causale automatique qui corrèle les anomalies à travers la pile applicative
complète. Davis est capable de remonter automatiquement la chaîne de causalité
d'un incident depuis le symptôme jusqu'à la cause racine.

\textbf{Points forts :} analyse causale la plus avancée du marché, autodécouverte
des topologies.

\textbf{Limites :} tarification extrêmement élevée (de l'ordre de 69~\$/mois
par \textit{full-stack host}) ; propriétaire et fermé ; ne convient pas aux petites
structures.

\subsection{Splunk + ITSI}

Splunk Enterprise est une plateforme d'analyse de journaux et de métriques
massivement utilisée en environnement d'entreprise. Son module \textit{IT Service
Intelligence} (\textsc{itsi}) \cite{splunk_itsi} ajoute des fonctionnalités de
détection d'anomalies basées sur le \textit{Machine Learning Toolkit} (\textsc{mltk})
et des corrélations d'événements. \caimp{} propose une intégration optionnelle
avec Splunk comme source de métriques secondaire, permettant aux organisations
déjà équipées de Splunk de bénéficier de l'IA explicative de \caimp{}.

\subsection{Alternatives open source : SigNoz, Netdata, Zabbix, Nagios}

\begin{description}[leftmargin=2cm]
  \item[SigNoz] Alternative open source à Datadog construite sur ClickHouse,
    supportant métriques, traces et journaux. Bonne interface, mais pas d'IA
    explicative.
  \item[Netdata] Monitoring temps réel à très faible latence, orienté serveur
    individuel, sans analyse multi-tenant.
  \item[Zabbix] Plateforme mature (depuis 2001) aux capacités étendues mais à
    la courbe d'apprentissage très abrupte ; pas d'IA native.
  \item[Nagios] Pionnier du monitoring (1999), aujourd'hui techniquement vieilli ;
    approche purement par seuils statiques.
\end{description}

\section{Tableau comparatif}

\begin{table}[H]
\centering
\caption{Comparaison des solutions de supervision selon 12 critères}
\label{tab:comparaison-solutions}
\small
\begin{tabularx}{\textwidth}{lXXXXXX}
\toprule
\textbf{Critère} & \textbf{Datadog} & \textbf{Grafana/Prom.} & \textbf{Dynatrace}
  & \textbf{SigNoz} & \textbf{Zabbix} & \textbf{\caimp{}} \\
\midrule
Prix (mois/hôte)  & 15--31~\$ & Gratuit    & 69~\$    & Gratuit     & Gratuit    & Gratuit \\
On-premise        & Non        & Oui        & Non      & Oui         & Oui        & Oui \\
IA native         & Partielle  & Non        & Avancée  & Non         & Non        & Oui \\
LLM explicatif    & Non        & Non        & Non      & Non         & Non        & \textbf{Oui} \\
RAG + runbooks    & Non        & Non        & Non      & Non         & Non        & \textbf{Oui} \\
Corrélation chang.& Non        & Non        & Partielle& Non         & Non        & \textbf{Oui} \\
Prévision 24h     & Oui        & Non        & Oui      & Non         & Non        & \textbf{Oui} \\
Dédup. alertes    & Oui        & Via AM     & Oui      & Partielle   & Non        & Oui \\
Multi-tenant      & Oui        & Limite     & Oui      & Oui         & Oui        & Oui \\
Déploiement \textless 10 min & Non & Non   & Non      & Non         & Non        & \textbf{Oui} \\
Dashboard answers-first & Non & Non        & Partielle& Non         & Non        & \textbf{Oui} \\
Tableau de bord   & Excellent  & Excellent  & Bon      & Correct     & Vieilli    & Moderne \\
\bottomrule
\end{tabularx}
\end{table}

\section{Techniques de détection d'anomalies}

\subsection{Détection par seuil statique}

La détection par seuil est la méthode la plus simple : une alerte est déclenchée
quand une valeur dépasse un seuil prédéfini. Sa simplicité est un atout, mais elle
présente deux défauts majeurs : elle génère des faux positifs pour les systèmes à
charge variable et ne s'adapte pas au comportement normal évolué d'un système.

\caimp{} implémente cette méthode comme premier filtre rapide, avec des seuils
configurables par variable d'environnement.

\subsection{Score Z et détection statistique}

Le Z-score mesure l'écart d'une observation à la moyenne en unités d'écart-type :

\begin{equation}
  z = \frac{x - \mu}{\sigma}
  \label{eq:zscore}
\end{equation}

où $\mu$ est la moyenne empirique calculée sur une fenêtre historique et $\sigma$
l'écart-type. Une valeur $|z| > 3$ est considérée comme anormale dans la
distribution normale standard (probabilité de 0,27~\%). Cette méthode s'adapte
automatiquement au comportement normal du système.

\subsection{Détection par taux de variation}

Cette méthode calcule la dérivée discrète de la série temporelle :

\begin{equation}
  \Delta_t = \frac{x_t - x_{t-1}}{x_{t-1}} \times 100\%
  \label{eq:rateofchange}
\end{equation}

Elle est particulièrement efficace pour détecter les saturations rapides (montée
soudaine du CPU lors d'un \textit{fork bomb}) que le Z-score pourrait manquer
sur une fenêtre historique longue.

\subsection{Forêt d'isolation (Isolation Forest)}

L'Isolation Forest \cite{liu2008isolation} est un algorithme non supervisé basé
sur la construction d'arbres de décision aléatoires. Son principe : les observations
aberrantes sont plus facilement \textit{isolées} par des partitions aléatoires.
Le score d'anomalie est inversement proportionnel à la profondeur moyenne
d'isolation. Bien que non implémenté dans \caimp{} v2, cet algorithme est
pressenti pour les versions futures.

\subsection{Lissage exponentiel de Holt-Winters}

La méthode de Holt-Winters \cite{holt1957forecasting, winters1960forecasting}
(ou double lissage exponentiel linéaire) modélise une série temporelle comme la
somme d'un niveau et d'une tendance linéaire :

\begin{align}
  L_t &= \alpha x_t + (1-\alpha)(L_{t-1} + T_{t-1}) \label{eq:hw_level} \\
  T_t &= \beta(L_t - L_{t-1}) + (1-\beta)T_{t-1}   \label{eq:hw_trend} \\
  \hat{x}_{t+h} &= L_t + h \cdot T_t               \label{eq:hw_forecast}
\end{align}

où :
\begin{itemize}
  \item $L_t$ est le niveau lissé au temps $t$ ;
  \item $T_t$ est la tendance lissée au temps $t$ ;
  \item $\alpha \in [0,1]$ est le paramètre de lissage du niveau ;
  \item $\beta \in [0,1]$ est le paramètre de lissage de la tendance ;
  \item $\hat{x}_{t+h}$ est la prévision à $h$ pas dans le futur.
\end{itemize}

Dans \caimp{}, les paramètres retenus sont $\alpha = 0.3$ et $\beta = 0.1$,
favorisant un lissage modéré qui réagit aux tendances sans amplifier le bruit.
Les intervalles de confiance sont calculés à $\pm 1.28 \sigma$ des résidus
(couverture approximative de 80~\%).

\section{Modèles de langage pour l'AIOps}

\subsection{Panorama des LLM ouverts}

L'émergence des modèles de langage à grande échelle (\textsc{llm}) ouverts a ouvert
une nouvelle dimension dans l'\textsc{aiops} : l'explication en langage naturel
des incidents. Parmi les modèles les plus pertinents pour une exécution locale :

\begin{description}[leftmargin=2cm]
  \item[Llama 3.1 8B] Le modèle retenu par \caimp{}. Publié par Meta AI en 2024
    \cite{llama3_meta2024}, il offre un excellent rapport qualité/performance en
    version quantifiée Q4\_K\_M : 5~Go de VRAM, latence inférieure à 2 secondes
    par requête sur du matériel courant.
  \item[Mistral 7B] Comparable à Llama 3 en termes de performance, avec une
    fenêtre de contexte étendue de 32K tokens.
  \item[Phi-3 Mini (3.8B)] Modèle compact de Microsoft, excellent pour l'inférence
    sur des machines à ressources limitées.
  \item[Gemma 7B] Modèle de Google, performant mais moins bien aligné pour les
    tâches opérationnelles.
\end{description}

\subsection{Quantification des modèles}

La quantification consiste à réduire la précision numérique des poids d'un modèle
(de float32 à int4, par exemple) pour diminuer son empreinte mémoire tout en
conservant une qualité d'inférence acceptable \cite{dettmers2022quantization}.
Le format \textsc{gguf} utilisé par Ollama permet la quantification
Q4\_K\_M (4 bits, méthode K-means) qui réduit la taille de Llama 3.1 8B de 16~Go
(float16) à environ 4.7~Go, rendant possible l'exécution sur un serveur standard.

\subsection{Le pattern RAG (Retrieval-Augmented Generation)}

Le \textsc{rag} \cite{lewis2020rag} est une technique d'augmentation du contexte
d'un \textsc{llm} par la récupération dynamique de documents pertinents depuis une
base de connaissances. Le schéma de fonctionnement est le suivant :

\begin{enumerate}
  \item L'\textbf{indexation} : les documents (runbooks, incidents passés) sont
    convertis en vecteurs d'embedding via un modèle dédié (\texttt{nomic-embed-text})
    et stockés dans une base vectorielle (pgvector \cite{pgvector2021}).
  \item La \textbf{recherche} : lors d'un incident, un vecteur est calculé pour
    le contexte de l'anomalie, et les $k$ documents les plus similaires sont
    récupérés par recherche de voisins approximatifs (cosine similarity).
  \item L'\textbf{augmentation} : les documents récupérés sont injectés dans le
    prompt du \textsc{llm} pour enrichir son contexte avant génération.
\end{enumerate}

Ce pattern permet à un \textsc{llm} à connaissance générale de s'adapter
spécifiquement au contexte d'une organisation, sans \textit{fine-tuning}
coûteux.

\section{Positionnement de CAIMP}

L'analyse de l'état de l'art fait apparaître un espace vacant entre les solutions
commerciales onéreuses (Datadog, Dynatrace) et les outils open source sans IA
(Grafana/Prometheus). Cet espace est celui des \textit{petites équipes} cherchant
une plateforme auto-hébergée, explicative et accessible.

\caimp{} se positionne sur cet espace avec les différenciants suivants :

\begin{itemize}
  \item \textbf{Explication en langage naturel} via un \textsc{llm} local (aucun
    outil open source ne propose cela) ;
  \item \textbf{Corrélation automatique des changements} (absente de toutes les
    solutions open source comparées) ;
  \item \textbf{Architecture événementielle moderne} avec NATS JetStream
    (plus légère que Kafka pour ce cas d'usage) ;
  \item \textbf{Déploiement en 10 minutes} sur n'importe quelle machine.
\end{itemize}

\section*{Conclusion}

L'état de l'art révèle que si les outils de monitoring ne manquent pas, aucun ne
combine simultanément l'auto-hébergement, l'accessibilité pour les non-experts
et l'explication en langage naturel par \textsc{llm} local. Cette lacune justifie
pleinement la création de \caimp{}. Le chapitre suivant traduit ces constats en
spécifications fonctionnelles et non fonctionnelles concrètes.
